\documentclass[a4paper,12pt]{article}
\usepackage[utf8]{inputenc} 
\usepackage[T1]{fontenc} 
\usepackage[ngerman,provide=*]{babel} 
\usepackage{lmodern}
\usepackage{geometry}
\usepackage{array}
\usepackage{booktabs}
\usepackage{hyperref}
\geometry{a4paper, margin=1in}

\begin{document}

\section*{1. Meilenstein: Ausarbeitung (bis 02.05.2025)}

\subsection*{Literaturrecherche und Einführung (2 Punkte)}

\textbf{Thema:} Entwicklung eines mehrschichtigen Sicherheitskonzepts für IP-Kameras mit Raspberry Pi in lokalen IoT-Netzwerken

\bigskip
\textbf{Ziel:}
\begin{itemize}
    \item Absicherung der Datenübertragung smarter IP-Kameras im lokalen Netz (zwei Raspberry Pis: Kamera \& Gateway).
    \item Vergleich von TLS, VPN (z. B. WireGuard) und SSH als Verschlüsselungsmethoden unter Praxisbedingungen.
    \item Einsatz von iptables-Firewall, MAC-Rotation und Authentifizierung zur Erhöhung der Gesamtsicherheit.
\end{itemize}

\bigskip
\textbf{Fragestellungen:}
{\renewcommand{\labelenumi}{\arabic*.}%
\begin{enumerate}
    \item Wie unterscheiden sich TLS, VPN und SSH in Bezug auf Sicherheit, Energieverbrauch und Umsetzung auf Raspberry Pi?
    \item Wie kann durch lokale Segmentierung (Namespaces) und Firewall-Regeln der Angriffsvektor reduziert werden?
    \item Welche Protokollkonfiguration ermöglicht eine effiziente, aber sichere Videoübertragung?
\end{enumerate}
}

\bigskip
\textbf{Recherchequellen:}
\begin{itemize}
    \item IEEE-Artikel zur IoT-Verschlüsselung und TLS-Leistung auf Embedded-Devices.
    \item Sicherheitsberichte zu IP-Kamera-Angriffen (z. B. Mirai-Botnetz).
    \item Dokumentationen zu VPN (OpenVPN, WireGuard) und SSH-Konfigurationen.
\end{itemize}

\bigskip
\hrule
\bigskip

\subsection*{Kritische Analyse einer Publikation (3 Punkte)}

\textbf{Ausgewählte Publikation:}\\
\textit{"Evaluation of TLS/SSL Implementations on Raspberry Pi for Secure IoT Communication"} von L. Lalande, Univ. Twente, 2024

\bigskip
\textbf{Analyse:}
\begin{itemize}
    \item \textbf{Stärken:}
    \begin{itemize}
        \item Praxisnahe Performance-Messung verschiedener TLS 1.3 Cipher-Suites.
        \item Session Resumption zeigt messbare Effizienzgewinne für häufige Verbindungen.
    \end{itemize}
    \item \textbf{Schwächen:}
    {\renewcommand{\labelenumi}{\arabic*.}%
    \begin{enumerate}
        \item Keine Berücksichtigung des Energieverbrauchs bei Dauerbetrieb.
        \item DTLS wurde getestet, aber ohne Erklärung für schlechte Ergebnisse.
    \end{enumerate}
    }
    \item \textbf{Verbesserungsvorschläge:}
    \begin{itemize}
        \item Ergänzung um Stromverbrauchs- und Langzeittests.
        \item Alternative Tunnelmethoden wie VPN (z. B. WireGuard) zum Vergleich implementieren.
    \end{itemize}
\end{itemize}

\bigskip
\textbf{Umsetzbare Ansätze für den Prototyp:}
\begin{enumerate}
    \item Zwei Raspberry Pis im LAN: einer als Kamera-Sender, einer als Gateway mit VPN-Server.
    \item Vergleich VPN-, TLS- und SSH-basierten Transports, Messung von Latenz, CPU-Auslastung und Sicherheit.
\end{enumerate}

\bigskip
\hrule
\bigskip

\subsection*{Projektidee und Konzeptskizze (4 Punkte)}

\textbf{Zeitplan (Beginn 13.04.2025, eine Woche früher!):}

\begin{center}
\begin{tabular}{@{}ll@{}}
\toprule
\textbf{Datum} & \textbf{Aufgabe} \\ \midrule
13.04.–19.04. & Literaturrecherche \& Einführung formulieren \\
20.04.–26.04. & Kritische Analyse & Verbesserungsvorschläge erarbeiten \\
27.04.–02.05. & Konzeptskizze verfassen \& DIN‑A4‑Seite finalisieren \\ \bottomrule
\end{tabular}
\end{center}

\bigskip
\textbf{Konzeptskizze:}
\begin{itemize}
    \item \textbf{Prototyp:}
    \begin{itemize}
        \item Pi A mit Kamera sendet Videodaten über TLS/WebSocket, VPN (z. B. WireGuard) oder SCP an Pi B.
        \item iptables auf Pi B begrenzt eingehende Verbindungen auf verschlüsselte Ports.
        \item MAC-Adresse des Pi A wird regelmäßig geändert.
    \end{itemize}
    \item \textbf{Technologien:}
    \begin{itemize}
        \item Raspberry Pi 4 + USB-Kamera
        \item TLS 1.3 via stunnel / Python
        \item SSH / SCP / SFTP
        \item iptables, Linux Namespaces, Cronjobs
    \end{itemize}
    \item \textbf{Begrenzungen:}
    \begin{itemize}
        \item Kein Cloud-Upload oder öffentliche Netzwerkverbindung
        \item Fokus liegt auf LAN-Szenario mit maximal zwei Geräten
    \end{itemize}
\end{itemize}

\bigskip
\textbf{Quellen (IEEE-Zitierweise):}
{\renewcommand{\labelenumi}{[\arabic*]}%
\begin{enumerate}
    \item L. Lalande, ``Evaluation of TLS/SSL Implementations on Raspberry Pi,'' B.Sc. Thesis, Univ. Twente, 2024.
    \item FasterCapital, ``Securing IoT Networks and Devices with Raspberry Pi,'' 2023. Online verfügbar: \url{https://fastercapital.com/topics/securing-iot-networks-and-devices-with-raspberry-pi.html}
    \item Cloudflare, ``What is the Mirai Botnet?,'' 2020. \url{https://www.cloudflare.com/learning/ddos/glossary/mirai-botnet/}
\end{enumerate}
}

\end{document}
